\chapter{The cost of a desert: a calibrated model and an engine}\label{ch:engineering}

Chapter~\ref{ch:foundations} reduced desert construction to a search: fix a
partial cover, walk the progression $\N=\N_0+kM$, and test the residual offsets.
This chapter answers the obvious next question---how long does that take?---and
then reports what happened when we ran it. The answer is a single exponent, and
the exponent is accurate enough that most of the rest of this thesis can proceed
by arithmetic rather than by experiment.

\section{The failure exponent}

Fix a partial cover with modulus $M=2\prod_{r\in\mathcal R}r$ and residual set
$U$ of prime offsets below $Q$. For a candidate $\N=\N_0+kM$, the covered
offsets are composite by construction; the candidate succeeds precisely when
\emph{every} residual complement $\N-q$, $q\in U$, is also composite.

Two effects set the probability. First, each complement is a number of size
about $\N$, so by the prime number theorem it is prime with probability about
$1/\log\N$. Second---and this is the roughness bonus of
Remark~\ref{rem:rough}---the cover has already forced $\N\equiv a_r$, so each
complement $\N-q$ avoids divisibility by every $r\in\mathcal R$ (for the
offsets that were \emph{not} killed by $r$), which raises its chance of being
prime. Mertens' theorem turns that into a multiplicative correction. Writing
\begin{equation}\label{eq:boost}
\mathrm{boost}\;=\;2\prod_{r\in\mathcal R}\frac{r}{r-1},
\end{equation}
the probability that one residual complement is prime is
$\approx\mathrm{boost}/\log\N$, and, treating the $|U|$ residual complements as
independent, the probability that a candidate survives all of them is
\[
\Bigl(1-\frac{\mathrm{boost}}{\log\N}\Bigr)^{|U|}\;\approx\;e^{-E},
\qquad
E\;=\;\frac{|U|\cdot \mathrm{boost}}{\log\N}.
\]
Thus $e^{E}$ candidates must be scanned per expected success. The independence
step is a heuristic---it is the Hardy--Littlewood picture applied to a fixed
arithmetic progression---and Chapter~\ref{ch:closure} shows that it is in a
precise sense the \emph{only} heuristic in play, because any correlation between
two complements is forced to act through small shared divisors, which the cover
already controls.

The formula explains the shape of the whole subject. Three levers move $E$:
\begin{itemize}
\item \textbf{More moduli} shrink $|U|$ but inflate $\mathrm{boost}$ and
$\log M$. The balance between these is not obvious, and \S\ref{sec:equilibrium}
shows it is an equilibrium rather than a budget constraint.
\item \textbf{A bigger $\N$} increases $\log\N$ and so lowers $E$ directly. This
is why the height game is easy and the budget game is hard: digits are the
cheapest resource in the problem.
\item \textbf{A higher target $Q$} raises $|U|$ roughly in proportion to
$\pi(Q)$, which is why record values of $\g$ get expensive fast.
\end{itemize}

\section{Calibration: does the model hold?}\label{sec:calibration}

A cost model that has never been checked is a wish. We therefore instrumented
every search to report, per slice of the progression, the measured pass rate
$\hat p$ of the probable-prime tests, the mean number of surviving offsets per
candidate, and the implied exponent
\[
\Ehat\;=\;\frac{(\text{mean survivors per candidate})\times\hat p}{1},
\]
which is the empirical analogue of $E$.

Across three independent covers and roughly $1.2\times10^{9}$ scanned
candidates, we find
\[
|\Ehat-E|\;\le\;0.1\ \text{nats},
\]
with per-slice values stable to within a couple of percent over hundreds of
slices. Two consequences matter. The obvious one is that projections in later
chapters can be made by arithmetic. The subtler one is a genuine bound on
hidden structure: if some unexploited arithmetic correlation made constructed
$\N$ systematically luckier than the model predicts, it would show up as a
persistent gap between $\Ehat$ and $E$. At $1.2\times10^{9}$ candidates, any
such effect is worth at most $0.1$ nats, a factor of $1.11$. We use this as a
quantitative constraint in Chapter~\ref{ch:closure}, where it rules out an
entire class of would-be mechanisms.

\section{The engine}\label{sec:engine}

The search is embarrassingly parallel and utterly dominated by one primitive: a
base-$2$ Fermat test on a number of the size of $\N$. Everything else is
bookkeeping designed to reduce the number of such tests.

\paragraph{Sieving.} For each residual offset $q$ and each small prime
$s\nmid M$, the candidates $k$ for which $s\mid\N-q$ form an arithmetic
progression in $k$, computable from one modular inverse. Marking those
progressions removes most $(k,q)$ pairs before any expensive test. The optimal
sieve bound is empirical: at $199$ digits we measure $B=10^{6}$ to beat
$B=3\cdot10^{6}$ ($901$ versus $816$ thousand candidates per second), because
beyond that the sieve costs more than the tests it saves.

\paragraph{Amortisation across variants.} A cover admits many
\emph{residue-swap variants}: reassigning some $a_r$ produces a different
progression with the same modulus $M$ and the same $|U|$, hence the same $E$.
These are free additional lottery tickets. Because they share $M$, they also
share every modular inverse in the sieve, so caching that table across variants
gave a measured $+47\%$ in throughput.

\paragraph{Survivor-ordered testing.} Within a block, candidates differ in how
many residual offsets survived the sieve. A candidate with $a$ survivors
succeeds with probability $(1-p)^{a}$, so hits concentrate sharply in the
low-survivor tail. Testing candidates in ascending order of $a$ is
result-preserving and measurably accelerates \emph{discovery} (a factor of
$1.83$ in expected time-to-first-hit at record scale).

\paragraph{Selective testing.} The natural next step is to \emph{skip} the
crowded candidates entirely. We implemented this with exact accounting: the
engine computes the survivor histogram, weights each candidate by
$(1-\hat p)^{a}$, and reports the fraction of total hit mass retained. This
keeps the model honest---$\hat p$ remains unbiased, so calibration still
works---while allowing a principled speed/coverage trade. At $199$ digits,
skipping the fullest $92\%$ of each block took the raw scan rate from
$1\,772$ to about $16\,300$ thousand candidates per second while retaining
$27.6\%$ of the hit mass: a net discovery rate about $3\times$ better.

The reason this works is worth stating, because it is not obvious. Each
candidate dies at its \emph{first} passing complement, so the expected number of
Fermat tests per candidate is about $1/\hat p\approx17$ regardless of how many
survivors it has. Testing cost is therefore nearly flat in $a$, while success
probability decays geometrically in $a$---so discarding high-$a$ candidates is
almost pure profit.

\paragraph{Hardware.} The same algorithm on GPU is a different animal. A CUDA
implementation using a bit-matrix sieve and Montgomery CIOS Fermat waves,
validated bit-for-bit against the reference implementation, sustains about
$2.5\times10^{6}$ tests/s on a T4 and $8.5\times10^{6}$ on an A100, versus
roughly $6\times10^{3}$ on four CPU cores. That ratio, about $430$, is
$\log 430\approx6.1$ nats---and the E-gap between the rungs the two platforms
can reach is likewise about $6$. The record frontier is logarithmic in compute,
and this is the cleanest illustration of it in our data: covers and strategy
were already saturated at CPU scale, and the GPU rung was bought purely with
throughput.

\section{Search strategy: depth, breadth, and the ratchet}\label{sec:ratchet}

Given a fixed compute budget, how should it be spent? We ran the experiment
both ways and the answer was unambiguous.

\emph{Breadth-first} designs scan many residue-swap variants shallowly. Because
the attainable $\N$ stays close to $\N_0$, every hit is small---good for the
threshold game---but the design has a hard ceiling: when the variant pool is
exhausted, the campaign ends with nothing.

\emph{Depth-first} designs scan one progression far. The runway is unbounded, and
because $\N$ grows only logarithmically with $k$, bad luck costs digits rather
than terminating the campaign. Our $150$-digit threshold record was found at
$k=6.21\times10^{10}$ of a full $[0,7.5\times10^{10})$ scan---at $63.8\%$
depth, a central draw against a cumulative expectation of $0.40$.

The synthesis is what we call the \textbf{ratchet}: design a cover, run it until
a hit, \emph{bank a certified record}, then re-tune the cover for the next rung
and repeat. At $Q\approx10^{5}$ this produced the descending ladder
$195\to179\to150$ digits, each rung a certified record in its own right. Under a
fixed budget the ratchet dominates a single grand-slam campaign for a structural
reason: it converts a lottery with one payout into a sequence of lotteries each
of which pays, and each rung also re-calibrates $\Ehat$ before more compute is
committed to the next.

\section{Results, including the ones that failed}\label{sec:campaigns}

Table~\ref{tab:records} listed the successes. Honesty requires the rest of the
distribution, so Table~\ref{tab:campaigns} reports the campaigns that produced
no record. In each case we give the survival probability---the model's own
prediction of how likely it was to see nothing, given the compute actually
spent---so the reader can judge whether the failure indicts the model or merely
the dice.

\begin{table}[h]
\centering\small
\caption{Campaigns that produced no record. ``Survival'' is $e^{-\lambda}$ where
$\lambda$ is the cumulative expected number of hits over the compute actually
spent; a value near $1$ means the campaign never had a real chance, a small
value means it was unlucky.}\label{tab:campaigns}
\begin{tabular}{lrrrl}
\toprule
Campaign & Target & Candidates & Survival & Outcome\\
\midrule
E10 (budget game, $Q=107\,720$) & $\g>107\,719$ & $2.4\cdot10^{8}$ & $0.17$ & unlucky, in-model\\
Campaign 2 ($Q=99\,991$, 168d cover) & $T(10^5)$ ladder & $5.6\cdot10^{8}$ & $0.22$ & unlucky, in-model\\
Campaign 3 ($Q=122\,000$) & $\g>119\,419$ & $4.0\cdot10^{8}$ & $0.93$ & stopped early\\
\bottomrule
\end{tabular}
\end{table}

Two of these were genuinely unlucky: a $0.17$ survival means the campaign had a
$83\%$ chance of producing a record and did not. That is a one-in-six outcome,
not a broken model, and in both cases the per-slice statistics tracked the
prediction throughout---which is precisely the discipline that lets us say so
with a straight face rather than reaching for an explanation. The third was
stopped by hand after $2.3\%$ of its budget when we learned that a parallel
effort was attacking the same target; its checkpoint is committed and it remains
resumable.

The methodological point is that a calibrated model turns a negative result into
information. Without $E$, three failed campaigns would be an embarrassment or an
excuse. With $E$, they are three consistent measurements of the same exponent,
and they are part of the $1.2\times10^{9}$-candidate calibration that
\S\ref{sec:calibration} rests on.

\section{Certification}\label{sec:certification}

Every record in this thesis is packaged so that its verification is independent
of its discovery. The pipeline is deliberately dull:

\begin{enumerate}
\item \textbf{Reconstruct.} From the committed cover and $(\N_0,M,k)$,
recompute $\N$ and check it against the stored decimal string.
\item \textbf{Exclude.} For every prime $q<\g$, produce a witness: either the
modulus $r$ with $r\mid\N-q$ (checked by division) or a strong base-$2$
compositeness witness (checked by modular exponentiation). Both are
unconditional; a prime cannot fail either test.
\item \textbf{Represent.} Verify that $\g$ is prime and that $\N-\g$ is prime,
the latter by an Atkin--Morain ECPP certificate~\cite{GoldwasserKilian,AtkinMorain}
generated by PARI/GP~\cite{PARI}.
\item \textbf{Re-verify independently.} A second checker, written against the
record file rather than the search code, re-derives every witness and validates
the certificate chain step by step---including the projective
elliptic-curve arithmetic---without calling the library that produced it.
\item \textbf{Hash.} A SHA-256 manifest fixes the artifacts.
\end{enumerate}

Step 4 is not ceremony. Re-verifying the two incumbent records at the start of
this programme is what established the baseline the rest of the work is measured
against, and the independent checker has since caught discrepancies in
intermediate artifacts that the primary pipeline reported as clean. Verification
is also far cheaper than discovery: it depends only on the final cover, the
finite residual check, and the certificate---never on reproducing the heuristic
path that found them.
